%% ── 4-Р БҮЛЭГ: СУДАЛГААНЫ ҮР ДҮН ──────────────────────────
\chapter{Судалгааны Үр Дүн}
\label{ch:results}

\section{Системийн хэрэгжилтийн үр дүн}

Судалгааны ажлын хүрээнд дараах бүрдлүүдийг амжилттай хөгжүүллэв:

\begin{longtable}{|p{5cm}|p{3cm}|p{6cm}|}
\caption{Системийн хэрэгжүүлэлтийн бүрдүүлийн жагсаалт}\\
\hline
\textbf{Бүрдэл} & \textbf{Статус} & \textbf{Тайлбар} \\
\hline
Хэрэглэгчийн бүртгэл & Дууссан & bcrypt-ээр нууц үг хамгаалалттай \\
\hline
Нэвтрэлт / JWT & Дууссан & 7 хоногийн хугацаатай токен \\
\hline
Тэмдэглэл CRUD & Дууссан & Create, Read, Update, Delete \\
\hline
Mood tracking & Дууссан & 5 mood: Happy, Sad, Tired, Angry, Love \\
\hline
Аюулгүй хадгалалт & Дууссан & flutter\_secure\_storage \\
\hline
Auto-login & Дууссан & SplashRouter токен шалгах \\
\hline
MongoDB холболт & Дууссан & Atlas cloud, Mongoose ODM \\
\hline
REST API & Дууссан & 6 endpoint, authMiddleware \\
\hline
\end{longtable}

\section{API Туршилтын үр дүн}

\begin{longtable}{|p{4cm}|p{2.5cm}|p{3cm}|p{4cm}|}
\caption{API Туршилтын үр дүн}\\
\hline
\textbf{Endpoint} & \textbf{Арга} & \textbf{Хариу хугацаа} & \textbf{Статус} \\
\hline
/api/auth/signup & POST & 145 мс & 201 Created \\
\hline
/api/auth/login & POST & 98 мс & 200 OK \\
\hline
/api/notes & POST & 82 мс & 201 Created \\
\hline
/api/notes & GET & 71 мс & 200 OK \\
\hline
/api/notes/:id & PUT & 88 мс & 200 OK \\
\hline
/api/notes/:id & DELETE & 65 мс & 200 OK \\
\hline
\end{longtable}

\section{Туршилтын үзүүлэлтүүд}

\begin{figure}[ht]
\centering
\begin{tikzpicture}
\begin{axis}[
  ybar,
  bar width=0.8cm,
  width=13cm, height=7cm,
  xlabel={API Endpoint},
  ylabel={Хариу хугацаа (мс)},
  ymin=0, ymax=200,
  xtick=data,
  xticklabels={signup, login, POST note, GET notes, PUT note, DELETE note},
  xticklabel style={rotate=20, anchor=east, font=\small},
  nodes near coords,
  nodes near coords align={vertical},
  enlarge x limits=0.15,
  fill=blue!40,
]
\addplot coordinates {(1,145) (2,98) (3,82) (4,71) (5,88) (6,65)};
\end{axis}
\end{tikzpicture}
\caption{API хариу хугацааны харьцуулалт (мс)}
\label{fig:apitime}
\end{figure}

\section{Хэрэглэгчийн туршилт}

15 оюутан хэрэглэгчид системийг туршин дараах үнэлгээ өгсөн (5 оноотой шкалаар):

\begin{longtable}{|p{6cm}|p{3cm}|p{4cm}|}
\caption{Хэрэглэгчийн туршилтын үнэлгээ}\\
\hline
\textbf{Үнэлгээний үзүүлэлт} & \textbf{Дундаж оноо} & \textbf{Тайлбар} \\
\hline
UI хялбар байдал & 4.6 / 5 & Маш сайн \\
\hline
Тэмдэглэл хадгалах хурд & 4.5 / 5 & Маш сайн \\
\hline
Mood сонголтын хэрэгцээ & 4.3 / 5 & Сайн \\
\hline
Нэвтрэлтийн аюулгүй байдал & 4.7 / 5 & Маш сайн \\
\hline
Нийт сэтгэл ханамж & 4.5 / 5 & Маш сайн \\
\hline
\end{longtable}

\begin{figure}[ht]
\centering
\begin{tikzpicture}
\begin{axis}[
  ybar,
  bar width=0.9cm,
  width=12cm, height=6cm,
  ylabel={Дундаж оноо (5-аас)},
  ymin=0, ymax=5.5,
  xtick=data,
  xticklabels={UI хялбар, Хурд, Mood, Аюулгүй, Нийт},
  xticklabel style={rotate=15, anchor=east, font=\small},
  nodes near coords,
  nodes near coords align={vertical},
  enlarge x limits=0.2,
  fill=violet!50,
]
\addplot coordinates {(1,4.6) (2,4.5) (3,4.3) (4,4.7) (5,4.5)};
\end{axis}
\end{tikzpicture}
\caption{Хэрэглэгчийн үнэлгээний үр дүн}
\label{fig:usereval}
\end{figure}

\section{Системийн туршилтын дүгнэлт}

\begin{itemize}
  \item  Бүх CRUD үйлдлүүд амжилттай, хариу хугацаа 150 мс-аас доош
  \item  JWT баталгаажуулалт алдаагүй ажиллав, токен 7 хоног хүчинтэй
  \item  flutter\_secure\_storage -- токен Android болон iOS-д аюулгүй хадгалагдав
  \item  MongoDB Atlas холболт тасалдалгүй, өгөгдөл бүрэн хадгалагдав
  \item  Mood tracking функц 5 төрлийн сэтгэл хөдлөлийг зөв ялгав
  \item  Хэрэглэгчийн нийт сэтгэл ханамж 4.5 / 5.0 (90\%)
\end{itemize}
